The preceding chapters traced tabular data understanding across two pillars (schema understanding and data understanding) and presented our contribution, Reflect-SQL~\cite{jain2027reflectsql}, to the former. This chapter sets out the research agenda that follows from that work. One concern unifies the entire agenda: making both pillars work at \emph{industry scale}, the regime of large, obsolescent schemas and tables too big for a model's context, where the methods that succeed on academic benchmarks break down. It is deliberately extensive and organized in two parts. The first part consolidates what the schema-understanding pillar has achieved and identifies the problems that remain open, since these directly extend our completed and ongoing work. The second part charts the direction we intend to pursue next: the data-understanding pillar, where the survey's open challenges (Section~\ref{sec:data}) define the frontier and where representing tabular data at scale, concretely as a graph, is the organizing commitment from which evidence retrieval and the remaining challenges follow.

%========================================================
\section{Future Directions in Schema Understanding}
\label{sec:future-schema}
The schema-understanding pillar has matured along the three axes identified in Section~\ref{sec:schema}. Obscure schemas are now routinely enriched with a semantic layer through knowledge-base and metadata construction; relevant schema elements are located through join-aware, iterative retrieval; and faithful queries are produced by multi-step pipelines that decompose, verify, and correct rather than generate in a single pass. Reflect-SQL advances this frontier by threading self-reflection through all three stages (knowledge-base construction, retrieval, and synthesis) and grounding it in the checkable signals that the SQL setting affords, reaching $72.03\%$ execution accuracy on BIRD. Yet several problems remain unsolved, and they define the near-term agenda.

\subsection{From static to continually evolving knowledge bases}
A recurring limitation across the knowledge-base literature, and a constraint our own work shares, is that the semantic layer is treated as a \emph{static} artifact: it is constructed once and thereafter fixed. Real schemas, however, drift: columns are added, business meanings shift, and each new query may reveal an undocumented facet of the data. The open problem is a \emph{continually evolving} knowledge base that updates itself from query logs, execution feedback, and user corrections, versioning schema semantics as they change rather than assuming a one-time build. Such a knowledge base would turn schema understanding from a preprocessing step into a lifelong-learning loop, in which every resolved query improves the system's model of the schema for the next one.

\subsection{Retrieval at industry scale}
Schema linking is accurate on academic benchmarks but becomes the binding constraint at industry scale, where a database may hold thousands of tables and tens of thousands of columns that cannot fit in a model's context window. The iterative, reflection-driven retrieval that Reflect-SQL employs is effective precisely because it reformulates the query when the retrieved context is judged insufficient, but every reformulation is another LLM call, so retrieval cost grows with the difficulty of the query. This mirrors a tension observed more broadly in structured retrieval, where coupling selection decisions to LLM invocations makes precision expensive~\cite{ye2023dater,wang2024chaintable}. A promising direction is to reframe schema retrieval from a query-time \emph{search} into addressing against \emph{pre-computed} structure: materializing foreign-key graphs, column-value indices, and description embeddings offline so that schema linking navigates directly to the relevant elements at near-constant cost, reserving LLM calls for genuine disambiguation rather than routine traversal. Decoupling retrieval from repeated model invocation is, in our view, the key to schema understanding that scales.

\subsection{Cross-database and multi-table join reasoning}
Reflect-SQL, like most current systems, is evaluated over a single database at a time. Enterprise text-to-SQL workflows are considerably harder: they span multiple databases and dialects, require long foreign-key chains and deeply nested queries, and demand reasoning about which of many near-duplicate tables is authoritative~\cite{lei2024spider2}. Extending self-reflective synthesis to this setting, where the space of join paths is combinatorial and a single wrong join silently corrupts the result, is an open problem. Join-aware retrieval that preserves relationship structure (Section~\ref{sec:schema}) is a necessary but not sufficient ingredient; the harder question is how reflection should reason about join \emph{correctness} when execution alone cannot distinguish a plausible-but-wrong join from the intended one.

\subsection{Reliable, reference-free verification}
Reflect-SQL's central bet is that reflection can be grounded in checkable signals (syntax, execution, and result--intent entailment scored by an LLM judge) so that queries can be validated without a gold reference. The reliability of that judgment is itself an open problem. LLM-as-a-judge is known to carry position, verbosity, and self-preference biases~\cite{zheng2023judging}, and a judge that is confidently wrong is more dangerous than one that abstains. Future work should pursue calibrated, reference-free verification: quantifying when a self-assessment can be trusted, learning when to escalate to a human or to additional evidence rather than accepting a flawed query, and stress-testing the semantic-validation dimensions against adversarial and near-miss queries where syntactic validity and semantic faithfulness diverge.

\subsection{The cost of reflection}
Reflection buys accuracy at the price of latency: iterative retrieval and iterative synthesis, each followed by multi-dimensional validation, multiply the number of LLM calls per query. For interactive analytics this cost is a deployment barrier. The open question is how to preserve the accuracy gains of reflection while collapsing its cost, for instance by distilling multi-step reflection into fewer passes, caching validated sub-queries and retrieval decisions across a session, and learning a stopping policy that halts refinement once marginal gains vanish. Efficiency here is not a secondary concern but a first-class research problem, and it foreshadows the efficiency--accuracy tension that dominates the data-understanding pillar below.

\subsection{Beyond SQL}
Finally, schema understanding is broader than SQL generation. Real analytical requests are frequently abstract and goal-oriented (``why did revenue dip last quarter?'') rather than a clean specification of a query, and answering them may call for dataframe or general code generation rather than a single SQL statement~\cite{he2024text2analysis,lei2024spider2}. Inferring the underlying analytical intent, and mapping it onto whichever execution substrate fits, generalizes the schema-linking problem and forms a natural bridge to the data-understanding pillar, to which we now turn.

%========================================================
\section{Setting a Direction for Data Understanding}
\label{sec:future-data}
Having established schema understanding through Reflect-SQL, our next research phase turns to the second pillar: data understanding, and concretely Table Question Answering, the problem of reasoning faithfully over the \emph{content} of tables. The survey literature frames this frontier through three challenges (Section~\ref{sec:data}): limited scope beyond retrieval and arithmetic, a lack of robustness on complex inputs, and limited generalization across representations. These challenges tell us \emph{what} is missing; the purpose of this section is to identify \emph{where} to begin. Our analysis points to a single organizing commitment from which the rest follow: \emph{how a large table is represented to the model}. The representations in use today are demonstrated almost exclusively on small, clean tables (Section~\ref{sec:representations}) and none scales to the hundreds of columns and thousands-to-millions of rows of real data; we therefore take \emph{representing tabular data at scale}, concretely as a graph, as the lead direction, and treat evidence retrieval, robustness, and higher-order reasoning as problems this representation is designed to unlock. We close by arguing that progress on all of them is gated by a benchmark that does not yet exist: Table Question Answering at industry scale.

\subsection{Representing large-scale tabular data as a graph}
\label{sec:future-representation}
The representations surveyed in Section~\ref{sec:representations} all flatten a table into a token stream (serialization), a pixel grid (image), or a fixed-width embedding (table encoders), and all are validated on tables small enough to fit whole in a model's input. None survives contact with industry-scale data: a table of hundreds of columns and millions of rows cannot be serialized within any context window, rendered legibly as an image, or embedded without discarding most of its content. The graph representation of Section~\ref{sec:graph} is, in our view, the most promising escape from this ceiling, but it has so far been demonstrated only on document graphs and small, complex tables. The direction we intend to pursue is to lift it to scale: to represent a large table as a graph in which cells, rows, columns, headers, and their groupings are typed nodes, and the table's relational topology (row--column membership, header spans, and foreign-key links across tables) forms the edges, materialized once and offline rather than reconstructed per query. Such a graph makes a table's structure explicit and \emph{addressable}: rather than asking a model to locate information inside a flattened dump, the structure itself records where everything is, so reaching a region of interest becomes a lookup rather than a scan. Making this representation hold at scale (how to construct the graph incrementally, index it for sub-linear access, and keep it current as the table evolves) is the central problem of this pillar, and the one on which the directions that follow depend.

\subsection{Evidence retrieval as addressing over the graph}
\label{sec:future-evidence}
Given such a representation, evidence retrieval (placing the handful of answer-bearing cells in front of the model before it reasons) changes character. Today it is the binding constraint: as tables grow to hundreds of rows and columns, an LLM must locate a few relevant cells amid overwhelmingly irrelevant content, and both context budget and accuracy degrade sharply as noise grows, the ``lost-in-the-middle'' effect~\cite{liu2024lostmiddle}. The deeper difficulty is that prevailing approaches couple cell selection to LLM invocation: full-table linearization performs no selection and exhausts the budget on large inputs; decomposition pipelines such as Chain-of-Table~\cite{wang2024chaintable} and DATER~\cite{ye2023dater} select rows and columns through a sequence of LLM calls, so every filtering decision costs an invocation and errors compound across stages; and even graph-based navigators such as GraphOTTER~\cite{li2025graphotter} traverse a structured table but score each hop with the LLM, tying latency to table size. In every case improving precision demands more LLM calls while reducing them risks discarding evidence, so an efficiency--accuracy tension appears fundamental. A pre-computed graph dissolves this tension by construction: because the groupings, statistics, and semantic partitions needed to navigate are query-independent, they can be computed offline, and retrieval becomes \emph{addressing} against that structure (navigating directly to the intersection of the relevant rows and columns) rather than an LLM-scored search, reserving the model for final synthesis. This decoupling of selection from reasoning is the direct content-side analogue of the schema-retrieval direction of Section~\ref{sec:future-schema}, and unifying the two under a single ``retrieval-as-addressing'' principle is the through-line we intend to develop across both pillars.

\subsection{Robustness on complex table structures}
The efficiency bottleneck is most acute precisely where the survey reports the largest human--model gaps: hierarchical, long, and multi-table inputs, where accuracies fall below $50\%$ against human performance around $83\%$ on benchmarks such as HiTab~\cite{cheng2022hitab} and MULTIHIERTT~\cite{zhao2022multihiertt}. Complex structure compounds the retrieval problem: multi-row headers carry span labels that flat serialization loses, and merged regions create implicit repetition that corrupts naive cell addressing. A robust system must therefore condition the table before reasoning (propagating hierarchical headers and normalizing merged cells) and carry that structure through to the representation, whether via specialized encoders in the tradition of TAPAS~\cite{herzig2020tapas} and TaBERT~\cite{yin2020tabert}, via a graph that encodes multi-row headers and merged regions as typed nodes and edges, or via structure-aware indices built over serialized views. Robustness to input complexity, rather than incremental gains on already-saturated simple benchmarks, is where we intend to concentrate.

\subsection{Reasoning beyond retrieval and arithmetic}
Once evidence can be delivered reliably, the harder question is what the model does with it. Many benchmarks are constructed by translating SQL or arithmetic expressions into natural language, so they reward matching synonyms to headers rather than genuine reasoning, and they are approaching saturation. The frontier is diagnostic, predictive, and prescriptive analysis (identifying insights, forecasting, and responding to abstract, goal-oriented requests) which requires understanding the table's overall context and the user's intent rather than surface retrieval~\cite{lei2024spider2,he2024text2analysis}. This is also where the two pillars converge: inferring analytical intent (a schema-understanding capability) and reasoning over content (a data-understanding capability) are two halves of the same problem, and a system that first understands structure and then reasons over content is the long-term target.

\subsection{Generalization across representations}
Finally, model performance varies by several percentage points with the input format, and systems transfer poorly across the diverse representations used by different benchmarks (Section~\ref{sec:representations})~\cite{sui2024table}. Promising remedies include standardizing serialization options, deriving serialized views from images of complex tables, and training on serialization-to-serialization tasks that convert one format into another to build robustness. Any representation we develop should be evaluated for representation-invariance from the outset, so that gains do not prove brittle to a change of input format.

\subsection{A large-scale benchmark for Table Question Answering}
A direction that conditions all of the above is measurement. The benchmarks that define the field (WikiTableQuestions, HiTab, MULTIHIERTT, and the others surveyed in Section~\ref{sec:data}) are built on tables of at most a few dozen rows and columns, and the retrieval-oriented ones are already approaching saturation; none exercises Table Question Answering at the scale where our representation and retrieval claims would actually be tested. There is, in short, no analogue of what BIRD~\cite{li2023bird} provided for Text-to-SQL: an industry-scale yardstick. Constructing one (tables with hundreds of columns and thousands-to-millions of rows, hierarchical and multi-table structure, realistic noise and obsolescent headers, and questions whose answers hinge on a vanishing fraction of the cells) is therefore a prerequisite rather than an afterthought. Such a benchmark would make the efficiency--accuracy trade-off measurable (cost per query against accuracy as table size grows), expose the brittleness that small, clean tables hide, and give the representation-at-scale and addressing directions above a concrete target against which progress can be judged. We regard building or curating this benchmark as part of the work, not a precondition to be supplied by others.

%========================================================
\section{Summary}
Taken together, these directions describe a coherent trajectory. In schema understanding, the immediate work is to make the semantic layer continually evolving, to scale retrieval by addressing rather than search, to extend faithful synthesis to cross-database joins, and to make reflection both reliable and cheap. In data understanding, our next pillar, the lead direction is to represent large tabular data at scale, concretely as a graph, so that evidence retrieval becomes addressing rather than search; robustness on complex structures, higher-order reasoning, and representation-invariance follow as targets this representation is meant to unlock, and a large-scale Table Question Answering benchmark is the instrument needed to measure progress on all of them. The unifying thread is that reliable tabular reasoning depends less on ever-larger models than on the architecture that delivers the right structure and the right evidence to them.
